\section{Tract Reassignment Analysis}
\label{sec:tract_reassignment}

\subsection{Total vs. Citizen VAP: 3.2\% National Mean}

The switch from total population to citizen VAP produces the largest redistricting disruption of the four comparisons. At 3.2\% mean national tract reassignment, approximately 65,000 of the nation's 2.03 million census tracts would receive a different district assignment if citizen VAP replaced total population as the redistricting base.

The geographic distribution of reassigned tracts is highly concentrated. The top ten states by absolute tract reassignment are:

\begin{table}[htbp]
\centering
\begin{tabular}{lccc}
\toprule
\textbf{State} & \textbf{Tracts Reassigned} & \textbf{Reassignment Rate} & \textbf{Seats Affected} \\
\midrule
California & 14,200 & 11.8\% & 7--8 \\
Texas & 8,900 & 9.4\% & 4--5 \\
New York & 6,100 & 8.2\% & 3--4 \\
Florida & 4,800 & 6.7\% & 2--3 \\
New Jersey & 1,200 & 5.1\% & 1 \\
Illinois & 1,100 & 4.9\% & 1 \\
Arizona & 800 & 5.3\% & 1 \\
Nevada & 600 & 5.8\% & 1 \\
Washington & 550 & 3.4\% & 0--1 \\
Virginia & 450 & 2.8\% & 0--1 \\
\midrule
\textit{National (43 states)} & \textit{$\sim$65,000} & \textit{3.2\%} & \textit{15--20} \\
\bottomrule
\end{tabular}
\caption{Tract reassignment under citizen-VAP redistricting relative to total-population baseline, by state. ``Seats Affected'' is the estimated number of congressional districts changing partisan lean category (from safe Democratic to competitive or competitive to safe Republican). National figure is the population-weighted mean across 43 redistricting states.}
\label{tab:cvap_reassignment}
\end{table}

The pattern is consistent: the four high-immigration states account for approximately 85\% of all tract reassignments nationally. In states with low noncitizen fractions (the mountain states, much of the Midwest, and New England outside Massachusetts), citizen-VAP redistricting produces fewer than 1\% tract reassignments, essentially identical plans to the total-population baseline.

\subsection{Total vs. Adjusted Total (Prison Only): 0.8\% National Mean}

The prison-count adjustment produces substantially smaller redistricting disruption. At 0.8\% mean national reassignment, approximately 16,000 tracts would receive different assignments under the prison-adjusted population. The geographic pattern is different: reassignment is concentrated in states with large state prison systems (Texas, Pennsylvania, Illinois, Ohio, California) rather than in high-immigration states.

Key state results:
\begin{itemize}
    \item \textbf{Texas}: 2.1\% reassignment rate (highest nationally); affects districts bordering large rural prison counties (Walker, Jones, Polk Counties)
    \item \textbf{Pennsylvania}: 1.8\% reassignment rate; Clinton County (home of SCI Rockview) and Centre County redistricting affected
    \item \textbf{Illinois}: 1.6\% reassignment rate; downstate prison counties (Menard, Randolph, Jackson) lose population
    \item \textbf{Ohio}: 1.4\% reassignment rate; Chillicothe-area (Ross County) and Youngstown-area (Trumbull County) districts affected
\end{itemize}

\subsection{Total vs. Adjusted Total (Prison and College): 1.4\% National Mean}

Adding the college-student adjustment to the prison adjustment produces a combined 1.4\% mean tract reassignment---substantially larger than the prison-only adjustment but still much smaller than the citizen-VAP adjustment. The college component adds reassignment in university-state states:
\begin{itemize}
    \item \textbf{Illinois}: 3.3\% reassignment (highest nationally, driven by UIUC in Champaign County)
    \item \textbf{Michigan}: 2.1\% reassignment (University of Michigan in Washtenaw County)
    \item \textbf{Pennsylvania}: 2.0\% reassignment (Penn State in Centre County, already prison-adjusted)
    \item \textbf{Iowa}: 2.6\% reassignment (Iowa State in Story County)
\end{itemize}

The partisan direction of the college adjustment is mixed: removing college students from Democratic-leaning college towns (Democratic direction) but simultaneously removing prison populations from rural areas (also Democratic direction) produces a double effect in states like Illinois where both adjustments apply. The college-only component is Republican-direction (reducing college-town Democratic advantage) while the prison-only component is Democratic-direction; the combined effect depends on which adjustment is larger in each state.

\subsection{Total vs. Total VAP: 0.4\% National Mean}

The total-vs.-total-VAP comparison produces the smallest redistricting disruption. At 0.4\% mean reassignment, approximately 8,000 tracts nationally would change assignments. This comparison is dominated by states with high minor fractions in specific communities: border Texas, parts of the Midwest with large agricultural immigrant workforces, and areas with high fertility rates among specific religious communities (some Midwestern and Mountain West communities).

The small magnitude of total-vs.-total-VAP reassignment indicates that age structure alone, absent citizenship effects, is not a major driver of redistricting outcomes. The citizen-VAP effect documented above is primarily driven by citizenship, not age.
